% !TeX program = xelatex
% !TeX encoding = UTF-8
\documentclass[12pt,a4paper]{article}
\usepackage[margin=2.5cm]{geometry}
\usepackage{amsmath,amssymb,amsthm}
\usepackage{xeCJK}
\setCJKmainfont{Microsoft JhengHei}
\usepackage{graphicx}
\usepackage[hidelinks]{hyperref}
\usepackage{fancyhdr}
\usepackage{needspace}

% Total physical pages; compile twice after page-count changes.
\makeatletter
\newcommand{\totalpages}{\@abspage@last}
\makeatother
\pagestyle{fancy}
\fancyhf{}
\renewcommand{\headrulewidth}{0pt}
\fancyfoot[C]{\thepage/\totalpages}
\fancypagestyle{plain}{%
  \fancyhf{}%
  \renewcommand{\headrulewidth}{0pt}%
  \fancyfoot[C]{\thepage/\totalpages}%
}
\setlength{\parindent}{0pt}
\setlength{\parskip}{0.5em}
\raggedbottom
\title{Linear System --- Homework 3}
\author{葉彥辰\quad P46154345}
\date{}

\begin{document}
\maketitle
\begin{center}Fall 2026\quad Due: Friday, Oct 16\end{center}

% 在各個 Solution 下撰寫解答；作答後可縮短或移除該處的 \vspace。
\section*{Exercise 1}
Obtain the transfer function of the following system
\[
\begin{aligned}
\dot x(t) &= -x(t)+2x(t-h)+u(t),\\
y(t) &= x(t).
\end{aligned}
\]

\textbf{Solution}
% 在這裡撰寫解答。

\begin{align*}
	&\begin{cases}
		\dot x(t) &= -x(t) + 2x(t - h) + u(t) \\
		y(t) &= x(t)
	\end{cases} \\
	\Rightarrow &\begin{cases}
		s\hat x(s) &= -\hat x(s) +2 e^{-sh} \hat x(s) + \hat u(s) \\
		\hat y(s) &= \hat x(s)
	\end{cases} \\
	\Rightarrow &\begin{cases}
		\hat x(s) &= {1\over s + 1 -2 e^{-sh} } \hat u(s) \\
		\hat y(s) &= \hat x(s) 
	\end{cases}
	\Rightarrow \hat y(s) =  {1\over s + 1 -2 e^{-sh} } \hat u(s) \\
\end{align*}
Hence, the transfer function
\[
\hat G(s) = {\hat y(s) \over \hat u(s)} = {1\over s + 1 - 2e^{-sh}}
\]

\newpage
\section*{Exercise 2}
Obtain a state space realization (controllable canonical form) of the transfer
function
\[
\hat G(s)=\frac{s^2+3s+2}{s^2+5s+6}.
\]
Is your realization minimal?

\textbf{Solution}
% 在這裡撰寫解答。

\begin{align*}
	\hat G(s) &= \frac{s^2+3s+2}{s^2+5s+6} \\ 
	& =  { (s + 1)(s + 2) \over (s + 2) (s + 3) }\\
	&=  { s + 1 \over s + 3} \\
	&= {-2 \over s+ 3} + 1\\
	\text {By} & \text{ Definition} \\
	&= C(sI - A)^{-1} B + D
\end{align*}
Hence
\[
A = -3, \quad B = 1, \quad C = -2, \quad D = 1
\]

\par\nopagebreak\vspace{3cm}

\Needspace{8cm}
\section*{Exercise 3}
Obtain a state space representation (controllable canonical form) of the following
transfer function:
\[
\hat G(s)=
\begin{pmatrix}
\dfrac{s^2+1}{s^2-1} & \dfrac{2}{s-1}
\end{pmatrix}.
\]

\textbf{Solution}
% 在這裡撰寫解答。

\[
\hat G(s) = \begin{bmatrix}
	\frac{s^2+1}{s^2-1} & \frac{2}{s-1}
\end{bmatrix} = \begin{bmatrix}
{ s^2 + 1 \over s^2 -1 } & {  2s + 2 \over s^2 -1 }
\end{bmatrix} = \begin{bmatrix}
{ 2  \over s^2 -1 } & {  2s + 2 \over s^2 -1 }
\end{bmatrix} + \begin{bmatrix}
1 & 0
\end{bmatrix} \Rightarrow \begin{cases}
	m &= 2 \\
	l&= 2
\end{cases} \Rightarrow n = m \times l = 4
\]

the deminsion of $I$ is $m \times m = 2 \times 2$, and we can know that $D = \begin{bmatrix}
	1 & 0
\end{bmatrix}$, directly.

\[
\hat G(s) = N(s) \Delta (s)^{-1}  + D\Rightarrow \Delta(s) = s^2 - 1 = \sum_{i = 0}^{l} \Delta_i s^i
\]
\[
N(s) = N_0 + N_1s = \begin{bmatrix}
	2 & 2
\end{bmatrix} + \begin{bmatrix}
0 & 2
\end{bmatrix}s
\]
Hence, 
\[
\Delta_0 = -I,\quad \Delta_1 = O, \quad \Delta_2 = I
\]
Finally, 
\[
A = \begin{bmatrix}
	O & I \\
	-\Delta_0 & -\Delta_1
\end{bmatrix} = \begin{bmatrix}
	0 & 0 & 1 & 0 \\
	0 & 0 & 0 & 1 \\
	1 & 0 & 0 & 0 \\
	0 & 1 & 0 & 0
\end{bmatrix}
\]
\[
B = \begin{bmatrix}
	O_{2\times2} \\ I_2
\end{bmatrix} = \begin{bmatrix}
	0 & 0 \\
	0 & 0 \\
	1 & 0 \\
	0 & 1
\end{bmatrix}
\]
\[
C = \begin{bmatrix}
	N_0 & N_1
\end{bmatrix} = \begin{bmatrix}
2 & 2 & 0 & 2
\end{bmatrix}
\]
\[
D = \begin{bmatrix}
	1 & 0
\end{bmatrix}
\]

--------------------------------


By
\[
\hat y(s) = \hat G(s) \hat u(s)\quad\Rightarrow \hat y(s) = {s^2 + 1 \over s^2 - 1} \hat u_1(s) + {2\over s - 1} u_2(s) = \left( 1 + {2\over s^2 - 1} \right) \hat u_1(s) + {2\over s - 1}\hat u_2(s)
\]
Let $ \hat q_1(s) = \displaystyle{\hat u_1(s) \over s^2 -1} $, and $ \hat q_2(s) = \displaystyle{\hat u_2(s) \over s - 1} $

then
\begin{align*}
	&\begin{cases}
		(s^2 - 1) \hat q_1(s) &= \hat u_1(s) \\
		(s -1) \hat q_2(s) &= \hat u_2(s)
	\end{cases} \\
	\Rightarrow &\begin{cases}
		\ddot q_1(t) - q_1(t) &= u_1(t) \\
		\dot q_2(t) - q_2(t) &= u_2(t)
	\end{cases}
\end{align*}
Let  $x = \begin{bmatrix}
	x_1 \\ x_2 \\ x_3
\end{bmatrix}$ with $x_1 = q_1$, $x_2 = \dot q_1$, and $x_3 = q_2$

\begin{align*}
	&\begin{cases}
		\dot x_2(t) - x_1(t) &= u_1(t) \\
		\dot x_3(t) - x_3(t) &=  u_2(t)
	\end{cases} \\
	\Rightarrow &\begin{cases}
		\dot x_1(t) &= x_2(t) \\
		\dot x_2(t) &= x_1(t) + u_1(t) \\
		\dot x_3(t) &= x_3(t) + u_2(t) \\
	\end{cases}
	\Rightarrow &\begin{bmatrix}
		\dot x_1 \\ \dot x_2 \\ \dot x_3 
	\end{bmatrix} = \begin{bmatrix}
	0 & 1 & 0 \\
	1 & 0 & 0 \\
	0 & 0 & 1 
	\end{bmatrix}\begin{bmatrix}
	x_1 \\ x_2 \\ x_3
	\end{bmatrix} + \begin{bmatrix}
	0 & 0\\
	1 & 0 \\
	0 & 1 	
	\end{bmatrix} \begin{bmatrix}
	u_1 \\ u_2
	\end{bmatrix}
\end{align*}
Hence, 
\[
A = \begin{bmatrix}
	0 & 1 & 0 \\
	1 & 0 & 0 \\
	0 & 0 & 1 
\end{bmatrix}, \quad B =\begin{bmatrix}
	0 & 0\\
	1 & 0 \\
	0 & 1
\end{bmatrix}
\]

By 
\begin{align*}
\hat y(s) &= \left( 1 + {2\over s^2 - 1} \right) \hat u_1(s) + {2\over s - 1}\hat u_2(s) \\
&= \hat u_1(s) + 2 {\hat u_1(s) \over s^2 -1} + 2 {\hat u_2(s) \over s - 1} \\
&= \hat u_1(s) + 2 \hat q_1(s) + 2 \hat q_2(s) \\
\Rightarrow y(t) &= 2 q_1(t) + 2 q_2(t) + u_1(t) \\
&= 2x_1(t) + 2 x_3(t) + u_1(t) \\
&= \begin{bmatrix}
	2 & 0 & 2
\end{bmatrix}\begin{bmatrix}
x_1 \\ x_2 \\ x_3 
\end{bmatrix} + \begin{bmatrix}
1 & 0
\end{bmatrix} \begin{bmatrix}
u_1 \\ u_2
\end{bmatrix}
\end{align*}




\
\section*{Exercise 4}
Obtain a state space realization (controllable canonical form) of the transfer
function
\[
\hat G(s)=
\begin{pmatrix}
\dfrac{1}{s-1} & \dfrac{1}{s-1} \\[8pt]
\dfrac{1}{s+1} & -\dfrac{1}{s+1}
\end{pmatrix}.
\]

\textbf{Solution}
% 在這裡撰寫解答。

\[
\hat G(s) = \begin{bmatrix}
	\displaystyle{ s+ 1 \over s^2 - 1} & \displaystyle{s + 1 \over s^2 - 1} \\
	\displaystyle{s - 1 \over s^2 - 1} & \displaystyle{-s + 1 \over s^2 -1}
\end{bmatrix}\Rightarrow \begin{cases}
m &= 2 \\
l &= 2 \\
n &= m \times l = 4
\end{cases}
\]
\[
\hat G(s) = N(s)\Delta(s)^{-1} + D(s), \quad D(s) = O_{2\times2}
\]
the deminsion of $I$ is $m \times m = 2 \times 2$
Hence, 
\[
\Delta(s) = \sum_{i = 0}^{l} \Delta_i s^i = s^2 - 1 \Rightarrow\begin{cases}
	\Delta_0 &= -I \\
	\Delta_1 &= 0 \\
	\Delta_2 &= I
\end{cases}
\]
\[
N(s) = N_0 + sN_1 = \begin{bmatrix}
	1 & 1 \\
	-1 & 1 
\end{bmatrix} + s\begin{bmatrix}
1 & 1 \\ 1 & -1
\end{bmatrix} 
\]

Finally, 
\[
A = \begin{bmatrix}
	O & I \\
	-\Delta_0 & -\Delta_1
\end{bmatrix} = \begin{bmatrix}
	0 & 0 & 1 & 0 \\
	0 & 0 & 0 & 1 \\
	1 & 0 & 0 & 0 \\
	0 & 1 & 0 & 0 
\end{bmatrix}
\]

\[B = \begin{bmatrix}
	O_{2\times2} \\ I_2
\end{bmatrix} = \begin{bmatrix}
	0 & 0 \\
	0 & 0 \\
	1 & 0 \\
	0 & 1
\end{bmatrix}
\]
\[
C = \begin{bmatrix}
	N_0 & N_1
\end{bmatrix} = \begin{bmatrix}
	1 & 1 & 1 & 1 \\
	-1 & 1 & 1 & -1
\end{bmatrix}
\]
\[
D = O
\]


\par\nopagebreak\vspace{3cm}

\Needspace{8cm}
\section*{Exercise 5}
\textbf{(a)} Obtain a state space realization (controllable canonical form) of the
following single-input single-output system:
\[
\ddot y-3\dot y-4y=\ddot u-2\dot u-8u.
\]

\textbf{Solution}
% 在這裡撰寫解答。

Under zero initial conditions, taking the Laplace transform gives
\[
(s^2-3s-4)\hat y(s)=(s^2-2s-8)\hat u(s).
\]
Hence,
\begin{align*}
\hat G(s)
&=\frac{\hat y(s)}{\hat u(s)}
 =\frac{s^2-2s-8}{s^2-3s-4}\\
&=\frac{(s-4)(s+2)}{(s-4)(s+1)}
 =\frac{s+2}{s+1}
 =1+\frac{1}{s+1}.
\end{align*}
Writing $\hat G(s)=N(s)\Delta(s)^{-1}+D$, we have
\[
\Delta(s)=s+1,\qquad N(s)=1,\qquad D=1.
\]
The reduced denominator has degree one. Its coefficient is $\alpha_0=1$,
and the numerator of the strictly proper part is $N_0=1$.
Thus, a first-order controllable canonical realization is
\[
\boxed{A=[-1],\qquad B=[1],\qquad C=[1],\qquad D=[1].}
\]
Equivalently,
\[
\boxed{\dot x(t)=-x(t)+u(t),\qquad y(t)=x(t)+u(t).}
\]
Indeed,
\[
C(sI-A)^{-1}B+D=\frac{1}{s+1}+1=\hat G(s).
\]

\Needspace{8cm}
\textbf{(b)} Is your realization minimal?

\textbf{Solution}

Yes. After pole-zero cancellation, the reduced transfer function is
\[
\hat G(s)=\frac{s+2}{s+1}.
\]
Its numerator and denominator have no common factor, and its denominator has
degree one. Hence, the minimum state dimension is one. Since the realization
in part (a) has exactly one state, it is minimal.

\Needspace{8cm}
\section*{Exercise 6}
Obtain a state space (controllable canonical form) realization of the following
input-output system:
\[
\begin{aligned}
\dot y_1+y_2 &= \dot u_2+u_1,\\
\dot y_2+y_1 &= \dot u_1+u_2.
\end{aligned}
\]

\textbf{Solution}
% 在這裡撰寫解答。
\par\nopagebreak\vspace{3cm}

\Needspace{10cm}
\section*{Exercise 7}
Determine whether or not the following system of linear equations has a solution.
If a solution exists, determine whether or not it is unique, and if not unique,
obtain an expression for all solutions.
\[
\begin{aligned}
x_1-x_2+2x_3+x_4 &= 5,\\
x_1-x_2+x_3 &= 3,\\
-2x_1+2x_2+2x_4 &= -2,\\
2x_1-2x_2-x_3-3x_4 &= 0.
\end{aligned}
\]

\textbf{Solution}
% 在這裡撰寫解答。
\par\nopagebreak\vspace{3cm}

\Needspace{14cm}
\section*{Exercise 8}
We consider here systems of linear equations of the form
\[
Ax=b
\]
with
\[
x=\begin{pmatrix}x_1\\x_2\\x_3\end{pmatrix}.
\]
For each of the following cases, we present the reduced row echelon form of the
augmented matrix $[A\ b]$. Determine whether or not the corresponding system of
linear equations has a solution. If a solution exists, determine whether or not
it is unique; if not unique, obtain an expression for all solutions and give two
solutions.

\textbf{(a)}
\[
\begin{pmatrix}
1 & 5 & 0 & 0\\
0 & 0 & 1 & 0\\
0 & 0 & 0 & 1
\end{pmatrix}
\]

\textbf{Solution}
% 在這裡撰寫解答。
\par\nopagebreak\vspace{3cm}

\Needspace{7cm}
\textbf{(b)}
\[
\begin{pmatrix}
1 & -2 & 0 & 6\\
0 & 0 & 1 & 4\\
0 & 0 & 0 & 0
\end{pmatrix}
\]

\textbf{Solution}
% 在這裡撰寫解答。
\par\nopagebreak\vspace{3cm}

\Needspace{7cm}
\textbf{(c)}
\[
\begin{pmatrix}
1 & 2 & 0 & 0\\
0 & 0 & 1 & 0\\
0 & 0 & 0 & 1\\
0 & 0 & 0 & 0
\end{pmatrix}
\]

\textbf{Solution}
% 在這裡撰寫解答。
\par\nopagebreak\vspace{3cm}

\Needspace{7cm}
\textbf{(d)}
\[
\begin{pmatrix}
1 & 0 & 0 & 0\\
0 & 1 & 0 & 0\\
0 & 0 & 1 & 0\\
0 & 0 & 0 & 1
\end{pmatrix}.
\]

\textbf{Solution}
% 在這裡撰寫解答。
\par\nopagebreak\vspace{3cm}

\Needspace{12cm}
\section*{Exercise 9}
Consider the three vectors
\[
\mathbf{v}_1=\begin{pmatrix}1\\2\\3\end{pmatrix},\quad
\mathbf{v}_2=\begin{pmatrix}2\\3\\4\end{pmatrix},\quad
\mathbf{v}_3=\begin{pmatrix}3\\4\\5\end{pmatrix}.
\]

\textbf{(a)} Express the vector
\[
\mathbf{v}=\begin{pmatrix}4\\5\\6\end{pmatrix}
\]
as a linear combination of $\mathbf{v}_1$, $\mathbf{v}_2$, and $\mathbf{v}_3$.

\textbf{Solution}
% 在這裡撰寫解答。
\par\nopagebreak\vspace{3cm}

\Needspace{8cm}
\textbf{(b)} Can every vector of the form
\[
\mathbf{v}=\begin{pmatrix}a\\b\\c\end{pmatrix}
\]
be expressed as a \textbf{linear combination} of $\mathbf{v}_1$, $\mathbf{v}_2$
and $\mathbf{v}_3$?

\textbf{Solution}
% 在這裡撰寫解答。
\par\nopagebreak\vspace{3cm}

\Needspace{4cm}
\textbf{(c)} Are the vectors $\mathbf{v}_1$, $\mathbf{v}_2$, $\mathbf{v}_3$
linearly independent?

\textbf{Solution}
% 在這裡撰寫解答。
\par\nopagebreak\vspace{2cm}

\Needspace{9cm}
\section*{Exercise 10}
\textbf{(a)} Find a basis for the \textbf{null space} of the matrix,
\[
A=\begin{pmatrix}
1 & 2 & 3 & 4\\
2 & 3 & 4 & 5\\
3 & 4 & 5 & 6
\end{pmatrix}.
\]

\textbf{Solution}
% 在這裡撰寫解答。
\par\nopagebreak\vspace{3cm}

\Needspace{4cm}
\textbf{(b)} What is the \textbf{nullity} of $A$?

\textbf{Solution}
% 在這裡撰寫解答。
\par\nopagebreak\vspace{2cm}

\Needspace{9cm}
\section*{Exercise 11}
\textbf{(a)} Find a basis for the \textbf{range} of the matrix,
\[
A=\begin{pmatrix}
1 & 2 & 3 & 4\\
2 & 3 & 4 & 5\\
3 & 4 & 5 & 6
\end{pmatrix}.
\]

\textbf{Solution}
% 在這裡撰寫解答。
\par\nopagebreak\vspace{3cm}

\Needspace{4cm}
\textbf{(b)} What is the \textbf{rank} of $A$?

\textbf{Solution}
% 在這裡撰寫解答。
\par\nopagebreak\vspace{2cm}

\Needspace{9cm}
\section*{Exercise 12}
Compute the \textbf{eigenvalues} and \textbf{eigenvectors} of the matrix,
\[
A=\begin{pmatrix}
2 & -3 & 0\\
2 & -3 & 0\\
3 & -5 & 1
\end{pmatrix}.
\]

\textbf{Solution}
% 在這裡撰寫解答。
\par\nopagebreak\vspace{3cm}

\Needspace{9cm}
\section*{Exercise 13}
Compute the \textbf{eigenvalues} and \textbf{eigenvectors} of the following matrix.
\[
A=\begin{pmatrix}
0 & 1 & 0 & 0\\
0 & 0 & 1 & 0\\
0 & 0 & 0 & 1\\
1 & 0 & 0 & 0
\end{pmatrix}
\]

\textbf{Solution}
% 在這裡撰寫解答。
\par\nopagebreak\vspace{3cm}

\end{document}
